% latexmk -xelatex --shell-escape -pvc memoria_tfg.tex

\documentclass [a4paper] {extarticle} % format
\usepackage [left=0.6in, right=0.7in, top=0.6in, bottom=0.6in] {geometry} % page margins

\usepackage {tikz}     % draw lines
\usepackage {hyperref} % hyperlinks
\usepackage {multicol} % columns
\usepackage {enumitem} % bullet lists configs
\usepackage {xcolor}   % colors
\usepackage[dvipsnames]{xcolor}

% icons
\usepackage {fontawesome5}
\usepackage {wasysym}
\usepackage {simpleicons}

\thispagestyle{empty} % no page number

%%%%%%%%%%%%%%%%%%%%%%
%%%%%%% NAME %%%%%%%%%
%%%%%%%%%%%%%%%%%%%%%%

\newcommand{\sectionlinethickness} {0.5pt}

% section size
\newcommand{\leftcolumwidth} {0.15}
\newcommand{\rightcolumwidth}{0.84}

\newcommand {\sectiontitle}[1] {
  \vspace*{-10pt}
  \begin {flushleft}
  \begin {minipage}[c]{\leftcolumwidth\textwidth}
    \begin {flushright}
      \!\MakeUppercase {#1}
      \hspace* {10pt}
    \end {flushright}
  \end {minipage}
  \begin{minipage}[c]{\rightcolumwidth\textwidth}
    \begin {tikzpicture}[remember picture, overlay,blend mode=multiply]
      \draw [line width=\sectionlinethickness, black] (0,0) -- (15,0);
    \end {tikzpicture}
  \end{minipage}
  \end {flushleft}
}

\definecolor{name}{HTML}{76C1C4}

\begin{document}

%%%%%%%%%%%%%%%%%%%%%%
%%%%%%% NOME %%%%%%%%%
%%%%%%%%%%%%%%%%%%%%%%

\begin {flushright}
\begin {minipage} [t] {\rightcolumwidth\textwidth}
\vspace {0.35\baselineskip}
\begin {minipage}[t]{0.45\textwidth}

\begin {tikzpicture}
\draw [line width=0.5pt, darkgray] (0,1) -- (0,1.5);
\end {tikzpicture}

\vspace {0.2cm}

\hspace {-0,39cm}
{\fontsize {20 pt} {20 pt}{\color{name} \bfseries \!\!\MakeUppercase {\hspace{11pt}iago}}}

\vspace{0.3cm}

{\fontsize {15 pt} {15 pt}{\!\MakeUppercase {Fernández Rego}}}

\vspace {0.2cm}

{\fontsize {15 pt} {15 pt}{\!\MakeUppercase {cientista da computaçom}}}

\vspace{0.2cm}

\begin {tikzpicture}
\draw [line width= 0.5pt, darkgray] (0,1) -- (0,1.5);
\end {tikzpicture}
\end {minipage}
\hfill
\vrule
\hspace*{5pt}
\begin {minipage}[t]{0.375\textwidth}
\begin {center}
\vspace{5pt}
\begin {tabular}{cl}
\color{gray}\faMapMarker & A Coruña, Galiza
\\[3.75pt] \color{gray}\faEnvelope  & \href {mailto:iago.fdez.rego@proton.me} {iago.fdez.rego@proton.me}
\\[3.75pt] \color{gray}\faGithub    & \href {https://github.com/rego-iago}{rego-iago}
\\[3.75pt] \color{gray}\faLink    & \href {http://rego.gal}{rego.gal}
%\faUniversity\faGlobe
\end {tabular}
\end {center}

\end {minipage}

\end {minipage}
\end {flushright}


%%%%%%%%%%%%%%%%%%%%%%
%%%%% SOBRE MIM %%%%%%
%%%%%%%%%%%%%%%%%%%%%%

\vspace*{-10pt}
\begin {flushleft}
\begin {minipage}[t]{\leftcolumwidth\textwidth}
  \begin {flushright}
  \!\MakeUppercase {sobre mim}
  \hspace* {10pt}
  \end {flushright}
\end {minipage}
\begin{minipage}[t]{\rightcolumwidth\textwidth}
  Graduado em Engenharia Informática de vinte e quatro anos. Gosto de resolver problemas e de entender o funcionamento interno de sistemas complexos ou concetos teóricos.\\
  Defendo umha tecnologia aberta, plural e ao serviço das pessoas.
\end{minipage}
\end {flushleft}

%%%%%%%%%%%%%%%%%%%%%%
%%%%% FORMAÇOM  %%%%%%
%%%%%%%%%%%%%%%%%%%%%%

\sectiontitle {formaçom}
\vspace{-15pt}

\begin {flushright}
\begin {minipage} [t] {\rightcolumwidth\textwidth}
  \begin {minipage} [t] {210pt}
    \href {https://estudos.udc.es/en/study/start/614G01V01} {\textbf {Grau em Engenharia Informática (inglês)}}\\
    \textbf {Universidade da Corunha}\\
    \hspace* {6pt} \textcolor{gray}{\textit{Setembro 2020 - Junho 2026}} \\ [4pt]
    Materias relevantes: Matemática Discreta, Algoritmos, Teoria da Computaçom, Aprendizagem Automática, Desenvolvemento de Sistemas Inteligentes,  Processamento de Linguagens.
  \end {minipage}
  \hfill
  \begin {minipage} [t] {200pt}
    \href {https://itg.es/} {\textbf {Desenvolvemento Software (full-stack)}}\\
    \textbf {Instituto Tecnológico de Galicia}\\
    \hspace* {6pt} \textcolor{gray}{\textit{Setembro 2023 - Junho 2024}} \\ [4pt]
    Trabalho colaborativo no desenvolvemento de umha ferramenta de autenticaçom e de umha aplicaçom de visualizaçom de dados, empregando JavaScript, Java e Angular
  \end {minipage}
\end {minipage}
\end {flushright}

%%%%%%%%%%%%%%%%%%%%%%
%%%%% PROJECTOS %%%%%%
%%%%%%%%%%%%%%%%%%%%%%

\sectiontitle {projetos}
\vspace{-15pt}

\begin {flushright}
\begin {minipage} [t] {\rightcolumwidth\textwidth}

  \href {https://github.com/yref-boop/ai-bias} {\textbf {LLM evaluation in regards to gender diversity and non-binary Galician language}}
  \hfill
  \vspace{1pt}
  \hrule
  \vspace {4pt} Projeto de fim de grau: trabalho de investigaçom no campo da socio-tecnologia que analisa o comportamento de LLMs comerciais ante a linguagem nom binária galega.
  \vspace{5pt}

  Também gosto de fazer pequenos projetos para testar ferramentas ou linguagens novas:
  \vspace{5pt}

  \begin {minipage} [t] {0.3\textwidth}
    \begin{minipage}[c]{100pt}
      \href {https://github.com/yref-boop/dotfiles} {\textbf {Dotfiles}}
    \end{minipage}
      \hfill
    \begin{minipage}[c]{8pt}
      \fontsize {8 pt}{8 pt}{\color{gray}\simpleicon{nixos}}
    \end{minipage}
    \vspace{1pt}
    \hrule
    \vspace {4pt} Configuraçons dos meus sistemas NixOS.
  \end {minipage}
  \hfill
  \begin {minipage} [t] {0.3\textwidth}
    \begin{minipage}[c]{100pt}
      \href {https://github.com/yref-boop/ann} {\textbf {Rede neuronal}}
    \end{minipage}
    \hfill
    \begin{minipage}[c]{9pt}
      \vspace{1pt}
      \fontsize {9 pt}{9 pt}{\color{gray}\simpleicon{haskell}}
    \end{minipage}
    \vspace{1pt}
    \hrule
    \vspace {4pt} Implementaçom de umha ANN do zero.
  \end {minipage}
  \hfill
  \begin {minipage} [t] {0.3\textwidth}
    \begin{minipage}[c]{100pt}
      \href {https://github.com/yref-boop/ann} {\textbf {Site web}}
    \end{minipage}
    \hfill
    \begin{minipage}[c]{9pt}
      \vspace{1pt}
      \fontsize {9 pt}{9 pt}{\color{gray}\simpleicon{gleam}}
    \end{minipage}
    \vspace{1pt}
    \hrule
    \vspace {4pt} Desenho e desenvolvemento do site web \href{rego.gal}{\underline{rego.gal}}.
  \end {minipage}
\end {minipage}
\end {flushright}

%%%%%%%%%%%%%%%%%%%%%%
%%%%% INTERESSES %%%%%
%%%%%%%%%%%%%%%%%%%%%%

\sectiontitle {Interesses}
\vspace{-15pt}

\begin {flushright}
\begin {minipage} [t] {\rightcolumwidth\textwidth}
  Sou uma pessoa moito curiosa, sempre aberta a investigar e aprender algo novo. Sendo a ciência da computaçom o meu foco principal, moitos dos meus interesses pessoais ficam nesse âmbito:

  \vspace {-8pt}
  \begin{multicols}{3}
    \begin {itemize} [noitemsep]
      \item Sócio-Tecnologia
      \item Programaçom Funcional
      \item Inteligência Artificial
      \item Avances Académicos
      \item FOSS
      \item Redes distribuida
    \end {itemize}
  \end{multicols}

  \vspace{-5pt}
  Também defendo a inter-disciplinidade como jeito de enriquecer qualquer área do conhecemento, tendo interesse pessoal por:

  \vspace{-8pt}
  \begin{multicols}{5}
    \begin {itemize} [noitemsep]
      \item Sociologia
      \item Linguística
      \item Mobilidade
      \item Urbanismo
      \item Desenho
    \end {itemize}
  \end{multicols}

\end {minipage}
\end {flushright}

%%%%%%%%%%%%%%%%%%%%%%
%%%%%% APTIDONS %%%%%%
%%%%%%%%%%%%%%%%%%%%%%

\sectiontitle {aptidons}
\vspace{-15pt}

\begin {flushright}
\begin {minipage} [t] {\rightcolumwidth\textwidth}

  \vspace {-8pt}
  \begin {itemize} [noitemsep]
    \item Ferramentas: JetBrains Suite, Neovim, Git, Bash, LaTeX, Nix, Inkscape.
    \item Linguagens de programaçom: C, Haskell, Ocaml, Python, Java, Gleam, Julia.
  \end {itemize}

  \textbf {Soft Skills} \\ [4pt]
  Para além do meu stack técnico, traio comigo as seguintes habilidades interpessoais:
  \vspace {-8pt}
  \begin{multicols}{3}
    \begin {itemize} [noitemsep]
      \item Adaptabilidade
      \item Colaboraçom
      \item Trabalho em equipe
      \item Resolución de problemas
      \item Independença
      \item Curiosidade
    \end {itemize}
  \end{multicols}

  \textbf{Línguas}: Espanhol \& Galego-português: \textcolor{gray}{\textit{nativo}}, Inglês: \textcolor{gray}{\textit{C1 (Cambridge)}}, Francês: \textcolor{gray}{\textit{B1 (DELF)}}

\end {minipage}
\end {flushright}

\end{document}
